\section{Base change in algebraic geometry}
\label{section-base-change}

\noindent
One motivation for the introduction of the language of schemes
is that it gives a very precise notion of what it means to
define a variety over a particular field. For example a variety
$X$ over $\mathbf{Q}$ is synonymous
(Varieties, Definition \ref{varieties-definition-variety})
with $X \to \Spec(\mathbf{Q})$ which is of finite type,
separated, irreducible and reduced\footnote{Of course algebraic
geometers still quibble over whether one should require $X$ to be
geometrically irreducible over $\mathbf{Q}$.}. In any case, the idea
is more generally to work with schemes over a given {\it base scheme},
often denoted $S$. We use the language: ``let $X$ be a scheme over $S$''
to mean simply that $X$ comes equipped with a morphism $X \to S$.
In diagrams we will try to picture the {\it structure morphism}
$X \to S$ as a downward arrow from $X$ to $S$.
We are often more interested in the properties of $X$ relative
to $S$ rather than the internal geometry of $X$.
For example, we would like to know things about the fibres of $X \to S$,
what happens to $X$ after base change, and so on.

\medskip\noindent
We introduce some of the language that is customarily used.
Of course this language is just a special case of thinking
about the category of objects over a given object in a category,
see Categories, Example \ref{categories-example-category-over-X}.

\begin{definition}
\label{definition-base-change}
Let $S$ be a scheme.
\begin{enumerate}
\item We say $X$ is a {\it scheme over $S$} to mean that $X$
comes equipped with a morphism of schemes $X \to S$.
The morphism $X \to S$ is sometimes called the
{\it structure morphism}.
\item If $R$ is a ring we say
$X$ is a {\it scheme over $R$} instead of
$X$ is a scheme over $\Spec(R)$.
\item A {\it morphism $f : X \to Y$ of schemes over $S$}
is a morphism of schemes such that the composition
$X \to Y \to S$ of $f$ with the structure morphism of $Y$ is
equal to the structure morphism of $X$.
\item We denote $\Mor_S(X, Y)$ the set of all morphisms
from $X$ to $Y$ over $S$.
\item Let $X$ be a scheme over $S$. Let $S' \to S$ be a
morphism of schemes. The {\it base change} of $X$
is the scheme $X_{S'} = S' \times_S X$ over $S'$.
\item Let $f : X \to Y$ be a morphism of schemes over $S$. Let $S' \to S$
be a morphism of schemes. The {\it base change} of $f$ is
the induced morphism $f' : X_{S'} \to Y_{S'}$ (namely the
morphism $\text{id}_{S'} \times_{\text{id}_S} f$).
\item Let $R$ be a ring. Let $X$ be a scheme over $R$.
Let $R \to R'$ be a ring map. The {\it base change} $X_{R'}$
is the scheme $\Spec(R') \times_{\Spec(R)} X$
over $R'$.
\end{enumerate}
\end{definition}

\noindent
Here is a typical result.

\begin{lemma}
\label{lemma-base-change-immersion}
Let $S$ be a scheme. Let $f : X \to Y$ be an
immersion (resp.\ closed immersion, resp. open immersion)
of schemes over $S$. Then any base change of $f$ is an
immersion (resp.\ closed immersion, resp. open immersion).
\end{lemma}

\begin{proof}
We can think of the base change of $f$ via the morphism
$S' \to S$ as the top left vertical arrow in the following
commutative diagram:
$$
\xymatrix{
X_{S'} \ar[r] \ar[d] & X \ar[d] \ar@/^4ex/[dd] \\
Y_{S'} \ar[r] \ar[d] & Y \ar[d] \\
S' \ar[r] & S
}
$$
The diagram implies $X_{S'} \cong Y_{S'} \times_Y X$,
and the lemma follows from Lemma \ref{lemma-fibre-product-immersion}.
\end{proof}

\noindent
In fact this type of result is so typical that there is a
piece of language to express it. Here it is.

\begin{definition}
\label{definition-preserved-by-base-change}
Properties and base change.
\begin{enumerate}
\item Let $\mathcal{P}$ be a property of schemes over a base.
We say that $\mathcal{P}$ is {\it preserved under arbitrary base change},
or simply that $\mathcal{P}$ is {\it preserved under base change}
if whenever $X/S$
has $\mathcal{P}$, any base change $X_{S'}/S'$ has $\mathcal{P}$.
\item Let $\mathcal{P}$ be a property of morphisms of schemes over a base.
We say that $\mathcal{P}$ is {\it preserved under arbitrary base change},
or simply that {\it preserved under base change} if whenever
$f : X \to Y$ over $S$ has $\mathcal{P}$, any base change
$f' : X_{S'} \to Y_{S'}$ over $S'$ has $\mathcal{P}$.
\end{enumerate}
\end{definition}

\noindent
At this point we can say that ``being a closed immersion'' is
preserved under arbitrary base change.

\begin{definition}
\label{definition-fibre}
Let $f : X \to S$ be a morphism of schemes.
Let $s \in S$ be a point.
The {\it scheme theoretic fibre $X_s$ of $f$ over $s$},
or simply the {\it fibre of $f$ over $s$},
is the scheme fitting in the following fibre product diagram
$$
\xymatrix{
X_s = \Spec(\kappa(s)) \times_S X \ar[r] \ar[d] &
X \ar[d] \\
\Spec(\kappa(s)) \ar[r] &
S
}
$$
We think of the fibre $X_s$ always as a scheme over $\kappa(s)$.
\end{definition}

\begin{lemma}
\label{lemma-fibre-topological}
Let $f : X \to S$ be a morphism of schemes.
Consider the diagrams
$$
\vcenter{
\xymatrix{
X_s \ar[r] \ar[d] & X \ar[d] \\
\Spec(\kappa(s)) \ar[r] & S
}
}
\quad\text{and}\quad
\vcenter{
\xymatrix{
\Spec(\mathcal{O}_{S, s}) \times_S X \ar[r] \ar[d] & X \ar[d] \\
\Spec(\mathcal{O}_{S, s}) \ar[r] & S
}
}
$$
In both cases these diagrams induce fibre squares of topological
spaces and in particular, the top horizontal arrow is a homeomorphism
onto its image.
\end{lemma}

\begin{proof}
Choose an open affine $U \subset S$ that contains $s$.
The bottom horizontal morphisms factor through $U$, see
Lemma \ref{lemma-morphism-from-spec-local-ring} for example.
Thus we may assume that $S$ is affine. If $X$ is also affine, then
the result follows from
Algebra, Remark \ref{algebra-remark-fundamental-diagram}.
In the general case the result follows by covering $X$ by open affines.
\end{proof}

\begin{lemma}
\label{lemma-local-ring-fibre}
Let $f : X \to S$ be a morphism of schemes. Let $x \in X$
with image $s \in S$. Thinking of $x$ as a point of $X_s$
(see Lemma \ref{lemma-fibre-topological}) we have
$$
\mathcal{O}_{X_s, x} \cong
\mathcal{O}_{X, x}/\mathfrak m_s\mathcal{O}_{X, x} \cong
\mathcal{O}_{X, x} \otimes_{\mathcal{O}_{S, s}} \kappa(s)
$$
\end{lemma}

\begin{proof}
As in the proof of Lemma \ref{lemma-fibre-topological} this reduces
to the case where $X$ and $S$ are affine. In the affine case the statement
translates to Algebra, Remark \ref{algebra-remark-local-ring-fibre}.
\end{proof}







